\chapter{Concluzii}
%\chapter*{Concluzii}
\label{conclusions}

Prezenta lucrare a abordat optimizarea sistemelor de recunoaștere a emoțiilor faciale bazate pe rețele neuronale cu impulsuri prin preprocesare adaptivă, urmărind atât validarea potențialului SNN-urilor ca extractoare de trăsături spațio-temporale, cât și integrarea practică a unui motor de inferență a emoțiilor într-o aplicație clinică funcțională.

\section*{Sinteza concluziilor experimentale}

Evaluarea pe setul de date CK+ prin \textit{10-fold cross-validation} a confirmat că modelul CSNN 3D antrenat nesupervizat prin STDP extrage trăsături temporale relevante pentru clasificarea emoțiilor faciale. Rezultatele comparative obținute sunt sintetizate după cum urmează:

\begin{itemize}
	\item \textbf{Modelul CSNN 3D (propus):} Atinge cea mai mare acuratețe medie (de \textbf{56.10\%}) folosinf \textit{Random Sampler}, depășind clasificatorul de referință (SVM pe pixeli bruți, 36.57\%) cu un câștig net de \textbf{19.53\%}. Acest fapt demonstrează că nucleele convoluționale specializate prin STDP codifică corelații spațio-temporale discriminative. Deși contraintuitiv, \textit{Random Sampler} excelează raportat la celelalte tipuri de eșantionare (\textit{bounding-box}, \textit{landmark-centrată}, bazată pe \textit{activitatea neuronală}) deoarece asigură o distribuție nepolarizată a spike-urilor, prevenind supra-ajustarea (\textit{overfitting}) algoritmului STDP pe regiuni hiper-specifice și păstrând intacte dinamici temporale subtile care altfel ar fi fost filtrate.
	\item \textbf{Modelul ANN echivalent:} A obținut \textbf{75.97\%} (utilizând, însă, fețe decupate), cuantificând costul utilizării unei reguli de învățare locale și nesupervizate față de optimizarea globală prin gradient (\textit{backpropagation}).
	\item \textbf{Arhitectura CNN-LSTM:} Proiectată ca referință de tip \textit{upper bound}, a atins \textbf{88.37\% $\pm$ 5.65\%}, confirmând superioritatea transferului de cunoștințe (ResNet18) și a modelării explicite a dinamicii temporale.
\end{itemize}

Analiza hărților de trăsături a evidențiat că modelul SNN funcționează ca un extractor de trăsături fiziologice: emoțiile cu activare musculară intensă (\textit{Happiness, Surprise, Disgust}) au obținut un \textit{recall} cel puțin de \textbf{83\%} pe primul fold, în timp ce emoțiile subtile (\textit{Contempt, Sadness}) au condus la confuzii sistematice.

Experimentele au subliniat un avantaj structural al arhitecturii SNN: filtrul On/Off și codificarea prin latență oferă o robustețe nativă la zgomotul vizual. Pixelii statici din fundal nu generează impulsuri, spre deosebire de modelele ANN unde fundalul poate scădea acuratețea la aproximativ 30--40\% în absența decupării. Cele patru metode de eșantionare propuse constituie baza unei analize experimentale detaliate, în urma căreia \textit{Random Sampler} a ieșit câștigător.

\section*{Sistemul Neuro-Guardian}

Aplicația dezvoltată demonstrează viabilitatea integrării unui motor de inferență al emoțiilor într-un flux clinic real. Elementele cheie ale implementării includ:
\begin{itemize}
	\item \textbf{Arhitectura hibridă:} WebSocket pentru latență redusă, TCP pentru comunicarea cu serverul din Python, PostgreSQL pentru date clinice și InfluxDB pentru monitorizarea seriilor temporale ale emoțiilor.
	\item \textbf{Mecanismul de histerezis:} Utilizarea pragurilor \textit{70/55} pentru alertele de tip \textit{Distress} și \textit{Pain} elimină oscilațiile la pragurile de decizie, oferind notificări stabile.
	\item \textbf{Validare complexă:} Combinarea testării \textit{white-box} (9 drumuri independente, Complexitate Ciclomatică = 9) cu integrarea \textit{top-down depth-first} realizată prin \textit{Jest}, \textit{Supertest} și \textit{Testcontainers} pentru izolarea dependențelor reale asigură o testare robustă a aplicației. 
\end{itemize}

\section*{Limitări și direcții viitoare de cercetare}

Lucrarea prezintă câteva limitări inerente:
\begin{enumerate}
	\item \textbf{Setul de date:} CK+ este colectat în condiții controlate, ceea ce poate limita generalizarea la emoții spontane.
	\item \textbf{Dezechilibrul claselor:} Numărul redus de exemple pentru \textit{Contempt} și \textit{Fear} afectează negativ performanța, STDP prioritizând trăsăturile frecvente.
	\item \textbf{Hardware:} Modelul rulează pe hardware convențional, beneficiile energetice ale SNN fiind maximizate doar pe cipuri precum Intel Loihi.
\end{enumerate}

Pe baza rezultatelor obținute, se propun următoarele direcții de extindere:
\begin{itemize}
	\item \textbf{Optimizarea altor hiperparametrilor:} Rafinarea pragului de activare ($\theta$) și a ratei de învățare ($\mu_{th}$) prin optimizare bayesiană.
	\item \textbf{Codificare temporală:} Înlocuirea \textit{latency coding} cu alte metode de codificare la nivelul întregii arhitecturi.
	\item \textbf{Antrenare End-to-End:} Implementarea unui strat \textit{Fully Connected Spiking} antrenat fie prin \textit{surrogate gradients}, fie prin STDP pentru a elimina necesitatea unui SVM extern.
	\item \textbf{Analiza robusteții:} Antrenarea și testarea arhitecturii și pe alte seturi de date mai complexe, cu mai mult zgomot de fundal.
\end{itemize}

În concluzie, lucrarea demonstrează că rețelele neuronale cu impulsuri sunt capabile să extragă trăsături relevante chiar și în regim nesupervizat, deschizând perspective reale pentru sisteme de recunoaștere a emoțiilor eficiente energetic, de tip \textit{Edge AI}.
